\subsection{Experimental Setup}
\subsubsection{Generating Counterfeit Examples}
We use three datasets: HumanEval, LeetCode, and ODEX \citep{wang2022execution}. HumanEval evaluates code generation of simple natural language descriptions, LeetCode is a harder dataset of the same flavor using programming interview practice problems, and ODEX tests knowledge of diverse Python libraries. To generate counterfeit examples, we first sample programs from CodeLlama (CL), DeepSeek-Coder Instruct (DS-I), and StarCoder (SC) at temperature $T=0.6$. Of the incorrect programs, we design a dataset-specific filter to remove incorrect programs that do not pass mild correctness criteria. For HumanEval, counterfeits are programs passing at least 10\% of EvalPlus \citep{liu2023your} tests. For LeetCode, counterfeits are programs that received a ``Wrong Answer'' verdict, which filters out programs that crashed during runtime or took too long to finish. For ODEX, counterfeits are programs that can be successfully parsed by \texttt{ast.parse} and are under 500 characters. In Fig. \ref{fig:sample-distribution}, we show the number of counterfeits generated by different models (left), benchmarks (middle), and problem difficulty levels (right), showing that counterfeits are widespread and occur in each setting.

\begin{figure}[ht!]
    \centering
    \begin{subfigure}[b]{0.45\columnwidth}
        \centering
        \includegraphics[width=\linewidth]{figures_v2/main/counterfeit_distribution/model_humaneval.pdf}
        % \subcaption{Model}
    \end{subfigure}
    \hfill
    \begin{subfigure}[b]{0.26\columnwidth}
        \centering
        \includegraphics[width=\linewidth]{figures_v2/main/counterfeit_distribution/benchmark_deepseek-instruct-33b.pdf}
        % \subcaption{Benchmark}
    \end{subfigure}
    \hfill
    \begin{subfigure}[b]{0.26\columnwidth}
        \centering        \includegraphics[width=\linewidth]{figures_v2/main/counterfeit_distribution/difficulty_livecodebench.pdf}
        % \subcaption{Difficulty}
    \end{subfigure}
    \caption{Distribution of correct, counterfeit, and incorrect samples by model, benchmark, and problem difficulty, showing their prevalence across all settings.}
    \label{fig:sample-distribution}
\end{figure}

\subsubsection{Code Understanding Tasks}
We evaluate counterfeits on the following tasks:

\textit{Correctness Checking}: The goal is to check whether a model-generated Python program (either correct or counterfeit) correctly implements a natural language (NL) specification. These specifications often include input-output examples. We use CoT with majority voting \citep{wei2022chain, wang2022self} and report accuracy.

\textit{Execution Prediction}: The goal is to predict the execution output of a given model-generated Python program on a specific input. We use an execution-based metric for correctness and report pass@1, the fraction of samples that are correct.

\textit{Repair}: The goal is to repair a given incorrect model-generated counterfeit program to correctly implement a given natural language specification. The model is not given any execution feedback other than the fact that the program is incorrect.

All three tasks are given to the language model (LM) in a few-shot setting. For GPT-3.5 and GPT-4, all tasks are prompted with chain-of-thought (CoT) \citep{wei2022chain}. For other models, we use CoT with majority voting ($N=10$) for correctness checking but not execution prediction, where we found it did not help. These tasks isolate semantic understanding of generated programs rather than merely asking for another unconditional code sample.

% A key difference from \citep{olausson2024repair} is that in our setup, the model is \textit{not given any information about which test failed}; it is only inform that the program is incorrect.
% Thus, the model has to localize and repair the issue without any external signal from the program interpreter.

\subsubsection{Dataset Creation}
Each set of samples is curated using a single dataset (such as HumanEval) and model (such as CodeLlama 34B). Each set is balanced and consists of 5 correct and 5 counterfeit programs for each problem (problems that do not have enough programs are discarded). Overall, across HumanEval, LeetCode, and ODEX, we generate $12$ different sets of samples, each consisting of $360$ to $1190$ programs. For correctness prediction, these datasets are used directly. For execution prediction, we randomly selected input-output examples, removing pairs that require complex arithmetic or execution. We also remove the problem statement and example input-output pairs so the model focuses on executing the code. For repair, we use the counterfeit samples in each set and discard the correct samples.
